% Terminal note of the homily: everything the speech itself must not carry ---
% the audience, the spoken word count and the pace it implies, the relation to
% the three readings of the expansive study, the route by which each quoted
% text reaches the page, and the exact loci --- followed by the References.
% Nothing here is read aloud.

\section*{Note on Sources and Delivery}
\runninghead{Sources and Delivery}

{\footnotesize

\textbf{Audience and occasion.} An adult parish assembly, at the Mass
\latin{Dominica vigesima post Pentecosten}, II classis, of the 1962
\work{Missale Romanum}, editio typica, pp.~412--413, nos.~1689--1698, for the
occurrence of 11~October 2026. The orations of the Maternity of the Blessed
Virgin Mary, commemorated on that date in a low or conventual Mass, belong to
another formulary and are not used.

\textbf{Length and pace.} The spoken body runs to 1,469 words, about 11.3 to
12.2 minutes at an unhurried 120 to 130 words a minute by arithmetic alone; no
one has delivered or timed it. It was read through in full, silently, for sense
and ease of speech.

\textbf{Relation to the reviewed readings.} The homily follows the first of the
full study's three readings of the whole Mass, \textit{The Word That Heals from
Afar}: Gregory's bodily presence sought from one absent nowhere and his healing
by command, Chrysostom's father healed with the son, Aquinas's faith growing on
the road, Augustine's due season and his ``fulfil Thy promise'', Chrysostom's
thanks ``even if they obtain it not'', the Secret's medicine, and the
Communion's unsung \latin{vivificavit me} beside \latin{Filius tuus vivit}. It
speaks the disagreement the study keeps on the ruler's first faith, Augustine's
``of no faith at all'' against Gregory's belief mixed with doubt. From the
second reading, \textit{The Song of the Exiles}, it takes Augustine's order of
prayer for the Introit, which the first and second readings share, the
Epistle's thanks for all things, and the Offertory's exiles with Augustine's
``whither your hope goeth before''. The third reading, \textit{The Dying Son
and the Two Peoples}, is left to the studies. The joins are the homily's own and
credited to no Father: the plan refused and the prayer granted, the father
learning the Introit's order, the Collect's secure mind as his on the road, the
Postcommunion's obedience as his going, and the exiles' hope as his road home.
So are the sentences on the altar, the Church's faith in the preacher's own
voice; no Father is said to join this Gospel to the Secret. The Alleluia is left
to the studies.

\textbf{Texts quoted, and their route.} Scripture is the Douay--Rheims,
Challoner's revision, at the canonical verses. The Communion is quoted in that
version of Ps 118:49 (119:49), which reads ``Be thou mindful'' for the Missal's
\latin{Memento} and lacks its added \latin{Domine}; the verse commended for
daily prayer is that exact Douay wording. The Introit's antiphon, the Collect,
the Secret and the Postcommunion are quoted in the anonymous English printed by
Eugene Cummiskey at Philadelphia in 1861, a historical study aid and not an
approved liturgical translation of the 1962 Missal; the antiphon gathers words
of Azarias's prayer (Dan 3) and is not a continuous verse. Augustine and
Chrysostom are quoted in the English of the Nicene and Post-Nicene Fathers.
Gregory and Aquinas are reported in English paraphrase of their Latin, and their
words are not quoted.

\textbf{Exact loci.} Introit, no.~1689; Collect, no.~1690; Epistle, Eph 5:20,
no.~1691; Gradual, Ps 144:15 (145:15), no.~1692; Gospel, Jn 4:48--50, 52--53,
no.~1694; Offertory, Ps 136:1 (137:1), no.~1695; Secret, no.~1696; Communion,
Ps 118:49--50 (119:49--50), no.~1697; Postcommunion, no.~1698. Chrysostom,
\work{Homilies on St John} 35, 2 (the father healed with the son) and 3
(thanks ``even if they obtain it not''). Augustine, \work{Tractates on John} 16,
3; \work{Enarrationes in Psalmos} 144, on vv.~15--17; 118, sermo 15, on v.~49;
136, 13. Gregory, \work{Homiliae in Evangelia} 28, 1. Aquinas, \work{Super
Ioannem} c.~4, lect.~7.

\par}

\Needspace{14\baselineskip}
\section*{References}
\runninghead{References}

{\footnotesize
\begin{itemize}\setlength{\itemsep}{0.15em}\setlength{\parskip}{0pt}
\item \work{Missale Romanum ex decreto Sacrosancti Concilii Tridentini restitutum, Summorum Pontificum cura recognitum}, editio typica (Typis Polyglottis Vaticanis, 1962), CMAA facsimile, pp.~412--413, nos.~1689--1698.
\item \work{The Roman Missal, Translated into the English Language for the Use of the Laity} (Philadelphia: Eugene Cummiskey, 1861), pp.~448--450 (XX Sunday after Pentecost, Introit, Collect, Secret, Postcommunion).
\item The Holy Bible, Douay--Rheims version, Challoner's revision: Ps 118:49--50 (119:49--50); 136:1 (137:1); 144:15 (145:15); Jn 4:48--53; Eph 5:20.
\item St John Chrysostom, \work{Homilies on St John} 35, Nicene and Post-Nicene Fathers, 1st ser., vol.~14 (New York, 1889), pp.~124--125.
\item St Augustine, \work{Tractates on the Gospel of John} 16, 3, Nicene and Post-Nicene Fathers, 1st ser., vol.~7. \work{Enarrationes in Psalmos} 118, sermo 15, on v.~49; 136, 13; and 144, on vv.~15--17 (Psalms CXIX, CXXXVII and CXLV in that translation), 1st ser., vol.~8.
\item St Gregory the Great, \work{Homiliae in Evangelia} 28, 1, Latin Wikisource transcription of Migne's text; Migne, PL 76, cols.~1210--1213; paraphrased, not quoted.
\item St Thomas Aquinas, \work{Super Evangelium S.~Ioannis lectura} c.~4, lect.~7, in \work{Opera}, editio altera Veneta, vol.~3 (Venice, 1745); paraphrased, not quoted.
\end{itemize}
\par}
